\subsection{The exact relative criterion and the missing existence assertion}

We retain the whole arithmetic group throughout. A cover is
geometrically connected, and ``constant'' refers to its deck group
over the displayed finite field. A $P$-split cover has a chosen rational
point over $P$; for a constant Galois group this makes its entire fibre
above $P$ rational.

\begin{claim}{B.CL1}{Proved}{The relative class-number test, with its group hypothesis}
Let $C/L$ be proper of genus at least two over a finite field of
characteristic $p$, and let $D\to C$ be a nontrivial constant Galois
\'etale $P$-split cover of degree $d$ prime to $p$. Put
$B=(\ker(J_D\to J_C))^0$. Then
\[
 \frac{\#J_D(L)}{\#J_C(L)}=\#B(L)\in\mathbf Z_{>0}.
\]
If the group has a nonidentity element of odd order, the condition
$B(L)[p]=0$ directly excludes a point-section collision at $P$.
This includes every odd prime cyclic cover used in B1.
The conclusion also holds for the remaining prime-to-$p$ groups
by the independently proved A.R1, but the odd-order stabilizer
argument by itself does not prove that broader group assertion.
\end{claim}
\begin{proof}
Pullback and norm satisfy $\pi_*\pi^*=[d]$. Hence
$J_C\times B\to J_D$ is an isogeny over $L$. The equality of
finite-field point counts under isogeny gives the asserted integer.
On $J_D(L)[p^\infty]$, multiplication by $d$ is invertible and
\[
 e=d^{-1}\pi^*\pi_*=d^{-1}\sum_{\delta\in G}\delta
\]
is an idempotent. Its complementary image lies in $B[p^\infty]$:
it is the image of the homomorphism $[d]-\pi^*\pi_*$, followed
by division by $d$ on this finite $p$-primary group. Thus
$B(L)[p]=0$ implies that all rational $p$-primary classes are
fixed by $G$.

Under a collision, A.R3 produces the unique lift $Q'$ above the
distinct point $Q$ with $T=[Q'-P']\ne0$ rational and $p$-primary.
If a nonidentity deck transformation fixes $T$, the effective
divisors $\delta Q'+P'$ and $Q'+\delta P'$ are distinct and
disjoint, and are linearly equivalent. They define a degree-two
pencil, so $D$ is hyperelliptic. Its hyperelliptic involution is
central. An odd-prime-order deck transformation descends to an
automorphism of the quotient projective line of that same order.
It has a geometric fixed point there; its permutation of the fibre,
which has at most two points, is trivial. It therefore has a fixed
point on $D$, contrary to the free deck action. This proves the
direct criterion. A.R1, proved from the configuration, boundary,
and cuspidal criteria in Part A, supplies the final broader conclusion.
The explicit even-order control in the norm-module discussion
explains why we have kept these two justifications distinct.
\end{proof}

\begin{claim}{B.CL2}{Proved}{Simultaneous character criterion in every genus}
Let $a:J_C(L)\twoheadrightarrow A$ be a nontrivial finite abelian
quotient of order prime to $p$, and use the Abel map based at $P$
to pull back its constant Lang quotient. Put
\[
 \Theta_{C,P,A}=\sum_{d=0}^{g-1}
       \sum_{E\in\operatorname{Div}^{+}_d(C)(L)}
                    [a([E-dP])]\quad\text{in }\mathbf F_p[A].
\]
Its relative class number is prime to $p$ if and only if
$\overline\psi(\Theta_{C,P,A})\ne0$ for every nontrivial
character $\overline\psi:A\to\overline{\mathbf F}_p^\times$.
For $A=\mathbf Z/\ell$ this is precisely B1's simultaneous
test; it is not enough to check one reduced character unless
all primitive characters form one Frobenius orbit.
\end{claim}
\begin{proof}
For a nontrivial complex character $\psi$ put
\[
 c_d(\psi)=\sum_{E\in\operatorname{Div}^{+}_d(C)(L)}
                 \psi(a([E-dP])),\qquad
 \kappa=[K_C-(2g-2)P].
\]
The Euler product, expanded by effective divisors, gives
$L_P(\psi,T)=\sum c_d(\psi)T^d$. A line bundle $M$ of
degree $d$ has $(q^{h^0(M)}-1)/(q-1)$ effective divisors in
its class. For $d>2g-2$, Riemann--Roch makes this number
independent of the class, so the nontrivial character sum
vanishes. For $g-1\le d\le2g-2$, apply Riemann--Roch to
$M$ and $K_C-M$; the class-independent term again vanishes.
This gives
\[
 c_d(\psi)=q^{d+1-g}\psi(a(\kappa))
                       c_{2g-2-d}(\psi^{-1}).
\]
In particular the coefficients with $d\ge g$ vanish modulo
every prime over $p$. This is the stated truncation at $g-1$.

On each closed fibre of the cover, its Frobenius translation
has the character eigenvalues $\psi(a_x)$ in the regular
representation. Factoring this permutation determinant and
multiplying the Euler factors proves
\[
 \frac{Z(D,T)}{Z(C,T)}=\prod_{\psi\ne1}L_P(\psi,T),
 \qquad
 \frac{\#J_D(L)}{\#J_C(L)}=
                       \prod_{\psi\ne1}L_P(\psi,1).
\]
Since $p\nmid|A|$, all characters remain distinct after
reduction in an algebraic closure of $\mathbf F_p$ and give
all reduced characters. The product is a unit exactly when
each factor is a unit. The preceding coefficient calculation
identifies the factors with the evaluations of $\Theta$.
If $A=\mathbf Z/\ell$, the rational factors are Galois
conjugates and their product is the cyclotomic norm. Over
$\mathbf F_p$, their orbits have length
$\operatorname{ord}_\ell(p)$, which proves the last assertion.
\end{proof}

\begin{claim}{B.CL3}{Proved}{The corrected pair estimate and its precise limitation}
For any quotient $A$ of $J_C(L)$ let $N=\#C(L)$ and let
$N_b$ be the occupancies of the Abel-image points in its cosets.
Then
\[
 \sum_bN_b(N_b-1)
 \le \frac{N^2}{|A|}
       +\frac{|A|-1}{|A|}(2g-2)^2q-N.                    \tag{B.CL.1}
\]
If all nonzero occupancies are multiples of $p$, the left side
is at least $(p-1)(N-N_0)$. In genus two, if $p$ generates
$(\mathbf Z/\ell)^\times$, a surjection with kernel
$H\subset J_C(L)$ is a class-number witness whenever
\[
                       N\ge2|H|,\qquad \ell>N.         \tag{B.CL.2}
\]
Neither $p>(2g-2)^2$ nor these estimates alone proves universal
existence of such a quotient. For $g>2$, occupancies of points
alone are not the character test in B.CL2.
\end{claim}
\begin{proof}
Character orthogonality gives
\[
 \sum_bN_b^2=\frac{1}{|A|}
       \left(N^2+\sum_{\psi\ne1}
          \left|\sum_{R\in C(L)}\psi(a([R-P]))\right|^2\right).
\]
Each nontrivial Lang character pulls back to a geometrically
nontrivial unramified rank-one local system: the Abel map
induces the isomorphism on geometric abelian fundamental groups.
Its $H^0$ and $H^2$ vanish, and its $H^1$ has dimension
$2g-2$ by the unramified Euler characteristic. It is a summand
of the cohomology of the finite \emph{\'etale} Lang pullback,
so its Frobenius eigenvalues have complex absolute value
$\sqrt q$, by the opened Weil theorem
\citep[Theorem~1.6 and formula~1.5.1]{WeilCurves}.
The trace formula bounds the inner sum by $(2g-2)\sqrt q$.
Subtracting $N$ proves (B.CL.1). For a positive occupancy
divisible by $p$, $N_b(N_b-1)\ge(p-1)N_b$; sum over $b\ne0$.

For $g=2$, B.CL2 reads $[0]+\sum_R[a([R-P])]$.
The primitive-root assumption makes $\Phi_\ell$ irreducible
over $\mathbf F_p$. Failure of the test therefore makes all
its coefficients equal. If $\ell>N$, an empty coset exists
and is nonzero, since $P$ occupies the zero coset. All nonzero
occupancies would consequently be divisible by $p$.
For a nonzero Jacobian class, at most two ordered pairs of
distinct geometric points have that difference: two different
representations give the unique degree-two pencil, and the
second pair must be $(\iota S,\iota R)$ when the first is
$(R,S)$. This proves the bound in all characteristics.
At most $2(|H|-1)$ points therefore share a coset with a
different point. Under (B.CL.2) at least two points occupy
singleton cosets, one of them nonzero. Its occupancy is not
divisible by $p$, a contradiction. Finally the unbounded
quotient-existence requirement and the effective-divisor terms
for $g>2$ have not been removed by this reasoning; the stated
limitations follow from the actual hypotheses just used.
\end{proof}

\begin{claim}{B.CL4}{Proved}{A finite all-constant-extension decision for each fixed prime}
Fix $(C,P)/\mathbf F_q$ and an odd prime $\ell\ne p$.
There is a finite procedure deciding whether B1 has an
order-$\ell$ witness over \emph{any} constant extension.
For each geometric nonzero character
$\chi:\pi_1(C_{\overline k})\to\mathbf F_\ell$, let $m$
be its least Frobenius period. Construct its unique $P$-split
constant cover over $k_m=\mathbf F_{q^m}$. It suffices to test
these finitely many covers at these first fields. If $A_\chi$
is Frobenius on their new geometric $p$-torsion, the exact
extension degrees preserving the witness are
\[
  n\ge1\quad\text{such that}\quad
                  \det(1-A_\chi^n)\ne0.                \tag{B.CL.3}
\]
A failed first-field test never becomes a witness at a later
field. A successful one remains successful under all
extensions of $p$-power degree.
\end{claim}
\begin{proof}
The character space is the finite vector space
$H^1_{\et}(C_{\overline k},\mathbf F_\ell)$ of dimension
$2g$, with its computable invertible Frobenius action.
Its nonzero elements, or their one-dimensional quotients,
are finite in number. An invariant character extends uniquely
to the semidirect arithmetic group by requiring it to be zero
on the chosen section $s_P$. This is exactly its geometrically
connected constant cyclic $P$-split cover. An extension permits
that same constant character exactly when its degree is a
multiple of $m$. Abelian class field theory, or the Lang
construction and the Abel isomorphism on geometric abelian
fundamental groups, identifies it with a quotient of
$J_C(k_m)$. These are the covers in B1, with the pointing fixed.

Pullback and norm split their $p$-primary Tate data, because
$p\ne\ell$. The quotient abelian variety's nonzero rational
$p$-torsion is therefore exactly the fixed space of $A_\chi$.
This proves (B.CL.3) and the first-field assertion: an existing
fixed vector remains fixed by every power. Conversely, a
constant cover over $k_{mn}$ with this geometric character and
split point is the base change of the unique split cover just
constructed, so no further twist has escaped the test.
For $n=p^a$, the polynomial identity
$X^{p^a}-1=(X-1)^{p^a}$ shows that $A_\chi^n-1$ is
invertible exactly when $A_\chi-1$ is invertible.

For an effective procedure, one may construct the finitely
many covers and compute their zeta numerators by point counts,
or compute the induced Frobenius on coherent cohomology and
its stable part by exact finite-field linear algebra. The
Artin--Schreier sequence identifies that stable part with
geometric \(p\)-torsion cohomology; nilpotent directions do
not affect $\det(1-F)$. Thus the procedure terminates for
every fixed $\ell$. It supplies no bound on the primes that
must be tested for the unrestricted existence assertion B1.
\end{proof}

\begin{claim}{B.CL5}{Open}{The remaining prescribed-point existence assertion}
For every prescribed $(C,P)$ as in B1, prove or refute the
existence of a prime $\ell$ for which at least one first-field
test in B.CL4 succeeds. This remains undecided here, including
the universal assertion for pairs $(p,g)$ with
$p>(2g-2)^2$. B.CL3 gives the corrected numerical sufficient
conditions; B.TW1--B.TW2 prove unpointed existence and quantify
the prescribed-point obstruction. The R8 example has the
explicit least-prime witness proved in B.W2--B.W3, and B.Q1
gives a plane-quartic instance. None of these removes the
arbitrary-curve and prescribed-point quantifiers.
\end{claim}
